\section{Introduction}
\label{sec:intro}

Wildlife management rests on reliable estimates of abundance. White-tailed deer are among the
most intensively managed ungulates in North America. Their density sets hunting quotas,
culling decisions and vehicle-collision mitigation, and the stakes extend beyond harvest.
Surveillance for chronic wasting disease is most efficient where infection risk is
highest~\cite{jennelle2018applying}. Deer-to-human transmission of SARS-CoV-2 was identified
in animals sampled during routine surveillance of exactly that
kind~\cite{pickering2022divergent}. Each of these decisions depends on knowing how many
animals are present.

For deer, one established way of estimating abundance is to drive transects at night with
spotlights or handheld thermal scopes. Road transects are practical but not neutral.
Spotlight surveys along roads carry a known bias~\cite{vercauteren2011managing}. A road used
as a transect line can also violate a central assumption of distance sampling, that animals
are distributed uniformly with respect to the line, because deer may be drawn to or kept
away from the road~\cite{diefenbach2025accounting}. Thermal imaging extends the survey into
darkness by detecting emitted heat, but it brings constraints of its own. Contrast falls as
the ground approaches body temperature, vegetation hides animals, and a distant deer
occupies only a few pixels~\cite{burke2019optimizing}. Any automated survey inherits all of
these at once.

Automating the survey is an obvious application of object detection, and the literature has
largely treated it as one. From camera traps~\cite{norouzzadeh2018automatically} to aerial
imagery~\cite{kellenberger2018detecting}, systems are trained on annotated animal imagery
and judged by per-image accuracy, typically mean average precision (mAP). The survey problem
is regarded as addressed once that number is high. This paper argues, and then measures,
that the two problems come apart. A detector is rewarded for finding animals in
\emph{frames}. A survey needs the number of distinct \emph{animals}. Between the two sits a
chain of decisions that no detection metric scores: associating detections into tracks,
deciding which tracks are real, and deciding which are duplicates of an animal already
counted.

That distinction fixes the yardstick we use throughout. A survey does not need a box placed
precisely on an animal; it needs the animal noticed once. We therefore score the detection
stage \emph{per animal}, over the whole of an animal's track, asking whether any frame
produced a box touching it. We report this alongside the conventional per-box metrics
rather than in place of them. Per-box mAP remains a valid measure of localization, but it
answers a question a survey does not ask, and read as a survey metric it can place the
difficulty in the wrong stage. Which stage limits a counting system is not something to
assume. In sonar fish counting, substituting a perfect detector recovers almost all of the
counting error, so there detection is the binding constraint~\cite{kay2022caltech}. A
thermal deer transect differs in ways that could move the constraint elsewhere. The animals
are faint, but each is visible over many consecutive frames as the vehicle passes, which
gives association and confirmation far more evidence to work with. Settling the question
needs an instrument that attributes error to each stage separately, applied to data
annotated at the level of individual animals.

We study this on a corpus built for the purpose: 32 FLIR thermal road transects from four
sites in Illinois, 521,930 frames at $640\times512$ and 60\,fps, with every individual deer
annotated as a distinct track. Because one annotated track is one animal, the corpus
supplies detection, tracking and count ground truth at once, which is the property the
decomposition depends on. Our central instrument is a three-stage decomposition of counting
error. For each ground-truth animal we ask three questions. Is it \reachedq, touched by at least one candidate track? Is it \primaryq, covered by some candidate better than that candidate covers any other animal, so that a one-count-per-track decision \emph{could} count it? And is it finally \countedq? The gaps between these three quantities attribute error to detection,
to association and to confirmation, and \cref{sec:decomposition} states the decomposition
formally.

With the loss localized, we test both remedies a practitioner would reach for. We enrich
the candidate pool by lowering the detector's confidence threshold, removing the tracker's
track-initialization gate, relaxing non-maximum suppression (NMS), and adding an appearance
term to association. We then replace the rule-based confirmation stage with learned
alternatives of increasing capacity, under three evaluation protocols. In each case the
question is the same: does the count improve?

The paper contributes three things. The first is the corpus: fully annotated thermal road
transects with per-individual identity, which makes detection, tracking and counting
evaluable against one set of annotations, released with the evaluation harness. The second
is the decomposition, together with a held-out protocol that evaluates counting only on
videos the detector never trained on. The third is a controlled test of where counting
error arises and whether improving detection or confirmation reduces it. The remaining
analyses support or qualify these three: the detector benchmark, the matching-criterion
analysis, the sensitivity of the confirmation rule, the limits of appearance matching at 27
pixels, and the transfer between sites.
